% =======================================================================================
\section{Limitations and Conclusion}
\label{sec:limitations}
% =======================================================================================

The simulator is a simplified \ac{lisa}: equal, constant arms, stationary Gaussian noise of known
\ac{psd}, no galactic foreground, no gaps or glitches, and exactly \RelatedWorkSourcesWord{}
detectable sources per window (\cref{app:forward-model-limits}). A real analysis needs the number
of sources inferred, the noise estimated with them, and windows stitched across their edges, as a
global fit does. The flow is checked on data from its own prior; out-of-distribution windows are
probed in \cref{app:evaluation-sampler} \textcolor{red}{TODO(new runs): outcome}. The sampler's
uniform grid can collapse the narrowest posteriors (\cref{fig:fm-sensitivity}), so widths are
read where they no longer change with the number of steps. Our validation tests are necessary
conditions; the Gaussian mock fails some of them, and only the reference chains give widths, for
a few windows. \textcolor{red}{TODO(new runs): where the flow stays wider than the
\acs{mcmc}, e.g.\ the loud sky and $\phi_-$.}

We have shown that conditional flow matching returns joint posteriors of overlapping \acp{gb}
under the full instrument response, across the \ac{lisa} band, in a fraction of a second per
window. The chart and the \ac{snr} ordering turn a posterior multimodal by symmetry into one a
flow on straight paths can learn, and the one-way network makes each draw cost only the source
stream. Validation needs tests in data space and known-answer mocks: without them, flows that
ignore the data pass.
